\begin{lemma}[extension of the invariant $(1,1)$ class]
\label{lem:bcpr-eta-extension}
\Verified\ The fibre class $\eta$ extends to a nonzero holomorphic section
\begin{equation}
 \boldsymbol\eta\in
 H^0\bigl(B,R^1f_*\Omega_X^1\bigr).
\label{eq:bcpr-eta-global-section}
\end{equation}
Consequently there is a class
$\mathfrak q\in H^{1,1}_{\bar\partial}(X)$ whose coherent-Leray edge
value is $\boldsymbol\eta$.  For every smooth fibre inclusion
$i_b\colon F_b\to X$,
\begin{equation}
 i_b^*\mathfrak q=[Q_F]_{\bar\partial}
 \in H^{1,1}_{\bar\partial}(F_b).
\label{eq:bcpr-q-restriction}
\end{equation}
\end{lemma}

\begin{proof}
Over the smooth marked locus, $\eta$ is flat, monodromy invariant, and of
type $(1,1)$ identically.  Equation
\Cref{eq:bcpr-eta-decomposition} identifies the corresponding holomorphic
section of the relative group as
\begin{equation}
 -\sigma_1\wedge u^{0,1}-6\sigma_2\wedge\gamma^{0,1}.
\label{eq:bcpr-relative-eta-section}
\end{equation}

At a finite filling write $t=s^M$, with $(M,a)=(3,2)$ or $(4,1)$.
The covering family is smooth at $s=0$, and
\cref{eq:bcpr-eta-decomposition} holds identically there, so
\cref{eq:bcpr-relative-eta-section} remains holomorphic.  It is invariant
under the cyclic action.  In equivariant frames $e_0,e_1$ for the relative
cotangent and $b_0,b_1$ for $R^1\mathcal O$, the actual invariant
$\mathbb C\{t\}$-basis from
\cref{prop:log-relative-direct-images} is
\begin{equation}
 e_0b_0,\quad s^ae_0b_1,\quad
 s^{M-a}e_1b_0,\quad e_1b_1.
\label{eq:bcpr-finite-invariant-basis}
\end{equation}
Exactness of cyclic invariants in characteristic zero therefore makes
\cref{eq:bcpr-relative-eta-section} a holomorphic descended section at
both finite points.

At the cusp, $\gamma$ and $u$ extend as sections of $R^1f_*\mathbb C$,
and their $(0,1)$ images form a basis of $H^1(W,\mathcal O_W)$ by
\cref{thm:structure-direct-images-audit}.  On the dense torus,
\[
 \sigma_i=d\zeta_i=(2\pi i)^{-1}d\log x_i.
\]
The classes of these logarithmic forms extend in the relative logarithmic
sheaf as
\begin{equation}
 \mathcal E|_{N_0}
 =\mathcal O_{N_0}[d\log x_1]\oplus
  \mathcal O_{N_0}[d\log x_2]
\label{eq:bcpr-cusp-log-basis}
\end{equation}
by \cref{prop:log-relative-direct-images}.  A cusp quotient shear changes
$d\log x_i$ only by a holomorphic multiple of $d\log t$, which is zero in
the relative logarithmic sheaf.  Thus
\cref{eq:bcpr-relative-eta-section} also extends holomorphically through
the cusp.  This concerns classes in the relative logarithmic sheaf; it
does not assert that the $d\log x_i$ are global ordinary one-forms on $X$.

The Gauss--Manin/Kodaira--Spencer map
\[
 \kappa_1\colon R^1f_*\mathcal E\longrightarrow
 \Omega_B^1(\log\Sigma)\otimes R^2f_*\mathcal O_X
\]
annihilates \cref{eq:bcpr-relative-eta-section} on the smooth locus: a
flat class which remains of type $(1,1)$ has zero $(0,2)$ Higgs component.
Both the section and $\kappa_1$ extend holomorphically, so the identity
theorem gives $\kappa_1(\boldsymbol\eta)=0$ on $B$.  The ordinary-cotangent
calculation in
\cref{prop:log-ks-divisor,prop:absolute-cotangent-direct-images,thm:h12-scalar}
identifies
\begin{equation}
 R^1f_*\Omega_X^1=\ker\kappa_1=:\mathcal K,
 \qquad
 \mathcal K\simeq\mathcal O\oplus\mathcal O(-1)^{\oplus2}.
\label{eq:bcpr-cotangent-kernel}
\end{equation}
This proves \cref{eq:bcpr-eta-global-section}.  Its value on every smooth
fibre is nonzero in Dolbeault cohomology, not merely as a differential form:
by \cref{eq:bcpr-eta-square} its square is the integral top class
$12\gamma\wedge u\wedge w\wedge\delta$, and therefore
\begin{equation}
 \int_{(F_b,\mathrm{or}_{\rm mark})}Q_F^2=12,
 \qquad
 \int_{(F_b,\mathrm{or}_{\mathbb C})}Q_F^2=-12
\qquad(b\in B^\circ),
\label{eq:bcpr-q-square-integral}
\end{equation}
where $\mathrm{or}_{\rm mark}$ is the orientation determined by
$(\gamma,u,w,\delta)$ and $\mathrm{or}_{\mathbb C}$ is the complex
orientation.  They are opposite because the determinant of the four marked
lattice columns in real coordinates is $y\Delta<0$.  Thus
$[Q_F]_{\bar\partial}\ne0$, so the holomorphic
section $\boldsymbol\eta$ is nonzero.  Explicitly, if
$Q_F=\bar\partial_F\alpha$ for a $(1,0)$-form $\alpha$, then
$Q_F^2=\bar\partial_F(\alpha\wedge Q_F)$; integration and Stokes' theorem
would contradict \cref{eq:bcpr-q-square-integral}.  Here
$\alpha\wedge Q_F$ has bidegree $(2,1)$, so its $\partial_F$ is zero in
complex dimension two and the displayed $\bar\partial_F$-boundary is also a
$d$-boundary.

Finally, coherent Leray over the curve gives the exact edge sequence
\begin{equation}
 0\longrightarrow H^1(B,f_*\Omega_X^1)
 \longrightarrow H^1(X,\Omega_X^1)
 \longrightarrow H^0(B,R^1f_*\Omega_X^1)
 \longrightarrow0,
\label{eq:bcpr-cotangent-edge}
\end{equation}
where $f_*\Omega_X^1\simeq\mathcal O(-2)$ by
\cref{prop:absolute-cotangent-direct-images}.  Choose $\mathfrak q$ over
$\boldsymbol\eta$.  Base change on a smooth fibre gives
\cref{eq:bcpr-q-restriction}.  Any two choices differ by the
one-dimensional left term in \cref{eq:bcpr-cotangent-edge}.  That term
restricts to zero on every smooth fibre.  The nonzero class
$d_1\mathfrak g$ also restricts to zero, because its restriction is
$[\partial_F\Gamma]_{\bar\partial}=0$; since
$H^1(B,\mathcal O(-2))$ is one-dimensional, the left term is exactly
$\operatorname{im}(d_1\colon H^{0,1}_{\bar\partial}\to
H^{1,1}_{\bar\partial})$.
\end{proof}
